% Main-text summary of recurring principles. Dense comparison: supplement/table_comparison_detail.tex.
\begin{table}[tbp]
\centering
\footnotesize
\singlespacing
\setlength{\tabcolsep}{5pt}
\renewcommand{\arraystretch}{1.3}
\caption{Recurring principles across laboratory and field rodent receivers. Each row states one point, the evidence that best illustrates it, and what that evidence does not establish. Fuller comparisons are in \stabref{tab:comparison-supp}{1}; study details are in \stabref{tab:evidence-supp}{2}.}
\label{tab:comparison}
\begin{tabular}{@{}>{\raggedright\arraybackslash}p{1.35in}>{\raggedright\arraybackslash}p{3.0in}>{\raggedright\arraybackslash}p{1.75in}@{}}
\toprule
\textbf{Principle} & \textbf{Illustrative evidence} & \textbf{What it does not show} \\
\midrule
Simple acoustic features often suffice
  & Tones approximating 50-kHz calls draw rats to the speaker, and frequency-matched tones reproduce most amygdala discrimination between call types \citep{WohrSchwarting2007,ParsanaEtAl2012a}.
  & That natural calls carry no further meaning. \\
Alarm calls carry different kinds of information
  & Marmot alarms grade risk by call rate; prairie-dog alarms specify escape suited to the predator \citep{BlumsteinArmitage1997,KiriazisSlobodchikoff2006}.
  & A single code shared across sciurids. \\
Experience gates action more than perception
  & Rats freeze to 22-kHz calls only after aversive experience, yet naïve rats still slow their heart and movement \citep{KimEtAl2010,OlszynskiEtAl2020}.
  & Whether predator calls are recognized without experience; tests conflict. \\
The body can respond without behavior
  & Spiny mice raise cortisol to owl calls without freezing or flight; young ground squirrels discriminate alarms by heart rate before they do so behaviorally \citep{EilamEtAl1999,Mateo1996}.
  & Which pathway produces the physiological change. \\
Parallel brain pathways allow these dissociations
  & Separate circuits support learned association, freezing, flight, and blood-pressure change \citep{LeDouxEtAl1988,FadokEtAl2017,TovoteEtAl2016,LiEtAl2021}.
  & That natural alarms use these circuits; most were mapped with tones or synthetic sounds. \\
Context selects the action
  & The same sound evokes flight or freezing depending on setting, and field receivers trade vigilance against foraging and refuge \citep{MongeauEtAl2003,BaackSwitzer2000}.
  & How costs are weighed in the brain. \\
Receivers track caller reliability
  & Ground squirrels discount unreliable callers, whereas marmots scrutinize them more closely \citep{HareAtkins2001,BlumsteinEtAl2004}.
  & Why the direction differs between species. \\
Declining responses have several causes
  & Responses recover when the caller changes, and irregular timing slows habituation to artificial sounds \citep{Hare1998,Davis1970,HerryEtAl2007}.
  & Which process underlies a given decline curve. \\
Meaningful sounds need not be threatening
  & Rats approach 50-kHz calls with a rise in accumbens dopamine, and oxytocin in auditory cortex enables mice to retrieve calling pups \citep{WilluhnEtAl2014,MarlinEtAl2015}.
  & Attraction to calls in free-living sciurids. \\
\bottomrule
\end{tabular}
\end{table}
